\section{Related Work}
\paragraph{Presentation generation.}
Early systems, including DOC2PPT~\citep{doc2ppt} and D2S/SciDuet~\citep{d2s}, treat document-to-slide generation as a summarization and composition problem. Recent work emphasizes language-driven planning and code generation: PPTAgent uses a two-stage editing approach~\citep{pptagent}, AutoPresent generates structured visuals through code~\citep{autopresent}, and DECKBench studies academic slide generation with iterative editing~\citep{deckbench}. These settings couple content planning with visual arrangement, but their evaluation units and input assumptions differ.

\paragraph{Presentation evaluation.}
Reference-based metrics compare an output with a target deck, as in parts of AutoPresent and DECKBench. Such comparisons do not directly establish compliance with an open-ended request that admits multiple valid decks. PPTEval~\citep{pptagent}, LecEval~\citep{leceval}, and PPTBench~\citep{pptbench} provide scoring criteria for content, design, or structure. PPTArena explicitly evaluates instruction-following and visual quality in editing tasks with ground-truth decks~\citep{pptarena}; SlideAudit provides a detailed design-defect taxonomy~\citep{slideaudit}; and SlidesGen-Bench evaluates content, aesthetics, and editability~\citep{slidesgen}. PresentBench uses manually designed, instance-specific binary checklists and evaluates multiple slide-generation systems~\citep{presentbench}. SlideGate instead constructs its task-specific specifications automatically from requests and parseable sources, merges and critiques candidates before freezing them, and adds programmatic references for objectively checkable instructions. Our AutoPresent comparison uses its reference-free design score because the tasks do not have unique target decks.

\paragraph{Model-based evaluation.}
G-Eval, MT-Bench, and Chatbot Arena connect model judgments to human evaluation~\citep{geval,mtbench,arena}. Instance-level checklist methods are particularly relevant: WildBench uses task-specific criteria~\citep{wildbench}, while CheckEval decomposes evaluation into Boolean questions~\citep{checkeval}. SlideGate applies this approach to presentations, combining frozen task specifications, source-grounded content checks, and page-level visual diagnosis. It reports dimensions and failure evidence rather than prescribing a single overall score.
